\documentclass[a4paper,11pt]{article}
\usepackage{exam-style-842-unit}
\examinehead{842 信息化概论 · 第 3 单元练习册}{练习一 参考答案与解析}
\begin{document}

\unitanswerscover{842 信息化概论 · 单元练习册}{第 3 单元参考答案}{练习一（计算机技术基础）的答案与解析；每道判断题均给出改正后的正确表述。}

\answerpart{一、单项选择题（10 题，每题 3 分，共 30 分）}

\q B。\begin{analysis}第一台电子计算机 ENIAC 于 1946 年 2 月 14 日在美国宾夕法尼亚大学诞生，是为计算导弹弹道而研制的（教材 2.1，第 16—17 页）。A 1936 年是图灵发表《理想计算机》论文的年份；C、D 分别是晶体管计算机和集成电路计算机发展之后的年份。\end{analysis}

\q C。\begin{analysis}冯·诺依曼提出作为一台电子计算机应有运算装置、存储装置、外设和输入输出装置，并采用存储程序原理，按这种思想设计的体系结构称为冯·诺依曼结构，现在所有的 PC 都用这个结构（教材 2.1，第 17 页）。A 图灵提出的是理想计算机（图灵机）模型；B 布尔建立的是逻辑代数；D 巴贝奇设计的是差分机和分析机。\end{analysis}

\q A。\begin{analysis}1936 年英国青年数学家图灵发表关于《理想计算机》的论文，严格描述了计算机的逻辑构造，提供了现代通用数字计算机的数学模型，并从理论上证明了它的可能性，为计算机诞生提供了理论准备（教材 2.1，第 16 页）。B 冯·诺依曼的贡献是存储程序与体系结构；C 布尔建立逻辑代数；D 艾肯研制了 Mark I。\end{analysis}

\q B。\begin{analysis}1854 年英国数学家乔治·布尔发表《思维规律研究》，建立了一套表示逻辑理论的符号及运算法则，把形式逻辑归结为代数演算，建立了逻辑代数（布尔代数），其运算规则与电路设计思想相同，是计算机逻辑设计的基础（教材 2.1，第 16 页）。A 莱布尼茨发明四则运算计算装置；C 帕斯卡制造齿轮式加法器；D 莫尔斯发明电报机与莫尔斯电码。\end{analysis}

\q A。\begin{analysis}第一代电子计算机 ENIAC 使用 1.8 万支电子管、1500 个继电器等，体积大、耗电多、速度慢，是第一代电子管计算机的代表（教材 2.1，第 17—18 页）。B 晶体管是第二代的主要器件；C 集成电路是第三代的主要器件；D 大规模集成电路是第四代的主要器件。\end{analysis}

\q A。\begin{analysis}第三代计算机最重要的特点是实现了“三化”——通用化、系列化、标准化，第三代计算机中还出现了小型计算机（教材 2.1，第 18 页）。B 中的微型化、网络化、智能化是后续发展的趋势；C 数字化、多媒体化、网络化是信息技术应用层面的特征；D 不是教材对第三代的概括。\end{analysis}

\q C。\begin{analysis}计算机部件技术主要有运算与控制技术、信息存储技术和信息输入输出技术，其中中央处理器处于核心地位（教材 2.2.3，第 21—22 页）。A 系统技术指系统结构、系统管理、系统维护和系统应用技术；B 器件技术指电子器件本身的技术；D 组装技术指制冷与提高组装密度等技术。\end{analysis}

\q A。\begin{analysis}计算机复杂逻辑的最基层线路为“与门”“或门”和“反相器”，由此组成的高一层线路有“组合逻辑”和“时序逻辑”两类（教材 2.2.2，第 20 页）。B 加法器、寄存器、计数器是部件；C 触发器、译码器、存储器也是由基本门电路构成的部件；D 组合逻辑与时序逻辑是门电路之上的逻辑类型，不是最基层线路。\end{analysis}

\q B。\begin{analysis}CAD（计算机辅助设计）就是利用计算机来进行产品设计，已广泛用于机械、船舶、飞机、大规模集成电路版图等设计，可提高设计质量、缩短设计周期、提高设计自动化水平（教材 2.3，第 22 页）。A CAM 是计算机辅助制造；C CAI 是计算机辅助教学；D 不是教材中的规范缩写。\end{analysis}

\q D。\begin{analysis}计算机技术作为信息技术的核心技术，具有“高技术”的特征：高度的创新性、高度的战略性、高度的增值性、高度的渗透性、高度的风险性（教材 2.4，第 23 页）。“高度的消耗性”不在其中，且与高技术“高增值、低消耗”的通行认识相反。\end{analysis}

\answerpart{二、多项选择题（4 题，每题 4 分，共 16 分）}

\q ABCD。\begin{analysis}操作系统的主要功能包括处理机管理、存储管理、文件管理、设备管理和作业管理（用户接口），四项全对。操作系统是最基本的系统软件，它管理计算机系统的硬件和软件资源，合理组织计算机工作流程，为用户提供方便的使用接口（大纲「计算机系统组成，计算机操作系统原理与功能」；教材 2.2.1 第 20 页指出计算机系统管理自动化是由操作系统实现的）。\end{analysis}

\q ABC。\begin{analysis}教材 2.2 把计算机技术粗分为系统技术、器件技术、部件技术和组装技术四类（第 19—24 页）。D 错误：网络技术是计算机技术与通信技术结合的产物，属于计算机技术的应用领域和外部环境，不是与系统、器件、部件、组装并列的技术分类。\end{analysis}

\q ABC。\begin{analysis}教材 2.5（第 24—25 页）指出二者的联系：计算机技术是信息技术的核心技术之一，计算机是信息接收、处理、存储的主要载体；二者都是高技术，都具有高技术的特征；计算机技术人才不能不懂 IT，IT 人才也离不开计算机，二者界限越来越模糊。D 错误：二者主要内容不同（计算机技术含硬件、软件及应用；信息技术强调对信息进行操作），发展历程和应用范围也不同。\end{analysis}

\q ABCD。\begin{analysis}教材 2.3（第 22—23 页）把计算机技术的应用概括为数值计算、研究方法、数据处理、计算机辅助设计、实时控制、模式识别、通信和图像文字处理、多媒体技术、网络技术与信息高速公路、教育等方面，四项全对。\end{analysis}

\answerpart{三、判断题（6 题，每题 2 分，共 12 分）}

\q √。\begin{analysis}第一台电子计算机 ENIAC 于 1946 年 2 月 14 日在美国宾夕法尼亚大学诞生，是为计算导弹弹道而研制的，比当时已有的计算装置快 1000 倍（教材 2.1，第 16—17 页）。\end{analysis}

\q √。\begin{analysis}冯·诺依曼提出电子计算机应有运算装置、存储装置、外设和输入输出装置，并采用存储程序原理，按此设计的体系结构称为冯·诺依曼结构，现在所有 PC 都采用这一结构（教材 2.1，第 17 页）。\end{analysis}

\q ×。改正：计算机发展的第四代以大规模集成电路为主要器件；晶体管是第二代计算机的主要器件，第四代还出现了微处理器和个人计算机。\begin{analysis}教材 2.1（第 18 页）指出：20 世纪 70 年代初期，随着大规模集成电路的出现，电子计算机进入第四代——大规模集成电路计算机，其成就是巨型机的发展以及微处理器与个人计算机的出现。\end{analysis}

\q √。\begin{analysis}操作系统是最基本的系统软件，负责管理处理机、存储器、文件、设备等硬件和软件资源，并作为用户与计算机之间的接口（大纲「计算机操作系统原理与功能」）；教材 2.2.1（第 20 页）也指出计算机系统管理自动化是由操作系统实现的，操作系统的基本目的在于最有效地利用计算机的软硬件资源。\end{analysis}

\q ×。改正：计算机技术与信息技术联系紧密，但主要内容、发展历程和应用范围均不相同，计算机技术是信息技术的核心子技术。\begin{analysis}教材 2.5（第 24—25 页）从主要内容、发展历程、应用范围三方面论证二者区别：信息技术先于计算机存在，外延比计算机技术大得多，计算机技术的应用范围只是信息技术应用范围的一部分。\end{analysis}

\q ×。改正：计算机技术包括硬件、软件及应用三方面的技术，并不只包括硬件技术。\begin{analysis}教材 2.5（第 24 页）指出：计算机技术主要包括计算机系统技术、器件技术、部件技术和组装技术，从内容上则包含硬件、软件及应用三方面的技术。\end{analysis}

\answerpart{四、填空题（6 空，每空 1 分，共 6 分）}

\q （1）ENIAC。\begin{analysis}1946 年 2 月 14 日诞生于美国宾夕法尼亚大学的第一台电子计算机名为 ENIAC（埃尼阿克），是为计算导弹弹道而研制的（教材 2.1，第 16—17 页）。\end{analysis}

\q （2）冯·诺依曼结构。\begin{analysis}现代计算机普遍采用冯·诺依曼结构，其核心思想是存储程序，由运算器、控制器、存储器、输入设备和输出设备五大部分组成（教材 2.1，第 17 页）。\end{analysis}

\q （3）组装技术。\begin{analysis}教材 2.2 把计算机技术分为系统技术、部件技术、器件技术和组装技术；组装技术与系统可靠性、维修调试方便性、生产工艺性和信息传递延迟密切相关（第 19—24 页）。\end{analysis}

\q （4）反相器。\begin{analysis}计算机复杂逻辑的最基层线路为“与门”“或门”和“反相器”（非门），由此组成组合逻辑和时序逻辑（教材 2.2.2，第 20 页）。\end{analysis}

\q （5）标准化。\begin{analysis}第三代计算机最重要的特点是实现“三化”——通用化、系列化、标准化，第三代还出现了小型计算机（教材 2.1，第 18 页）。\end{analysis}

\q （6）实时控制。\begin{analysis}实时控制也叫过程控制，就是用计算机对连续工作的控制对象实行自动控制，要求及时采集信号、计算处理并发出调节信号（教材 2.3，第 22 页）。\end{analysis}

\answerpart{五、简答题（3 题，每题 8 分，共 24 分）}

\q 参考答案要点：
（1）第一代（1946 年至 20 世纪 50 年代末）：主要器件是电子管，代表是 ENIAC。体积大、功耗高、速度慢、可靠性差、价格昂贵，主要用于军事和科学计算，用机器语言编程。
（2）第二代（1960 年起）：主要器件是晶体管。引入快速磁芯存储器和磁鼓、磁带、磁盘等外存储器，运算速度由电子管的几千次每秒提高到几十万次每秒，重量、体积、功耗和售价大幅下降。
（3）第三代：主要器件是集成电路。最重要的特点是实现“三化”——通用化、系列化、标准化，第三代计算机中还出现了小型计算机。
（4）第四代（20 世纪 70 年代初期起）：主要器件是大规模集成电路。一是巨型机发展，在核武器研制、导弹及航空航天、气象预报、卫星图像处理、经济预测等领域作用重大；二是微处理器与个人计算机出现，计算机在世界范围内大规模普及。
（5）规律总结：器件变革是计算机划代的标志，计算机发展是由应用和成本互相驱动、由数学和物理学共同推动的过程，目前计算机体系结构仍是冯·诺依曼结构。（教材 2.1，第 16—19 页）
\begin{analysis}得分点：四代的器件与时间各 1 分共 4 分，各代主要特点共 3 分，规律总结 1 分。常见失分：把“三化”错记到第四代；漏答第四代的两个成就（巨型机、微处理器与个人计算机）；只写器件不写特点。本题对应 2023 年真题“计算机基本结构分几个阶段及特点”。\end{analysis}

\q 参考答案要点：
（1）计算机技术可粗分为系统技术、器件技术、部件技术和组装技术四类（教材 2.2，第 19—24 页）。
（2）系统技术：把计算机作为一个完整系统所运用的技术，包括系统结构技术（使系统获得良好解题效率和合理性能价格比）、系统管理技术（由操作系统实现，提高资源利用率和吞吐能力）、系统维护技术（功能检查程序和自动诊断程序）、系统应用技术（程序设计自动化和软件工程）。
（3）器件技术：电子器件是计算机系统的物质基础，最基层线路为与门、或门、反相器；器件从真空电子器件到超大规模集成电路的变革是计算机划代的标志，使组装密度提高约 4 个数量级、速度提高 5—6 个数量级、功耗和价格大幅降低。
（4）部件技术：主要有运算与控制技术（中央处理器处于核心地位，指令重叠、指令并行、流水线、超高速缓冲存储器、微程序技术）、信息存储技术（形成超高速缓存—主存—外存的层次结构，半导体存储器是主存的主要手段）、信息输入输出技术（穿孔卡片纸带、软盘、光文字阅读机、语音图像输入；击打式与非击打式打印机、显示终端）。
（5）组装技术：与系统可靠性、维修调试方便性、生产工艺性和信息传递延迟密切相关，核心是制冷（直接制冷、水冷、强制风冷，多采用风冷）和提高组装密度，随计算机换代不断向小型化、微型化发展。
\begin{analysis}得分点：四个分类名称 2 分，每类主要内容 1 分共 4 分，器件划代或存储层次等关键点 2 分。常见失分：把网络技术列入四类技术；只写分类名称不写内容；混淆“器件技术”与“部件技术”。教材 2.2 第 19—24 页。\end{analysis}

\q 参考答案要点：
（1）联系（教材 2.5，第 24 页）：
（1）计算机技术是信息技术的核心技术之一，计算机是信息接收、处理、存储的主要载体；
（2）计算机技术和信息技术都是一项高技术，都具有高技术的一些特征；
（3）计算机技术人才不能不懂 IT，而 IT 人才更是离不开计算机，二者界限变得越来越模糊。
（2）区别（教材 2.5，第 24—25 页）：
（1）主要内容不同：计算机技术主要包含硬件、软件及应用三方面的技术，以离散数学、算法论、语言理论、控制论、信息论、自动机理论等为基础；信息技术主要强调对“信息”进行操作的技术，主要支柱是通信技术、计算机技术和控制技术，即“3C”技术；
（2）发展历程不同：信息技术处理的是信息，没有计算机之前也存在信息技术，如电报、电话，甚至烽火传信、飞鸽传书都属于信息技术；计算机从 ENIAC 至今不到 100 年，即使从帕斯卡加法器算起也不到 400 年；
（3）应用范围不同：计算机技术是信息技术的子技术，信息技术外延极大，还包括电视、录音、光盘、网络、数码相机、扫描仪等，计算机技术的应用范围只是信息技术应用范围的一部分。
（3）结论：二者联系紧密但不能互相替代，计算机技术是信息技术的核心子技术，同时又以计算机这一实体为研究对象，有自己的发展规律。
\begin{analysis}得分点：联系三点 3 分，区别三点 5 分。常见失分：只说“二者是一回事”或只罗列联系不答区别；发展历程部分漏掉“没有计算机之前已存在信息技术”这一关键论据。教材 2.5 第 24—25 页。\end{analysis}

\answerpart{六、综合题（1 题，共 12 分）}

\q 参考答案要点：
（1）主要发展方向（4 分）：教材 2.6.1 指出计算机技术的主要发展方向是生物计算机、光子计算机、量子计算机和纳米计算机。
（1）生物计算机：以蛋白质分子、DNA 核酶等生物分子作为转换开关和存储介质，并行处理能力强，能耗只有电子计算机的十亿分之一，运算速度可比最新一代计算机快 10 万倍；受生物技术水平限制，预计 10—20 年后进入实用阶段。
（2）光子计算机：以光子代替电子、光互连代替导线互连、光硬件代替电子硬件、光运算代替电运算；并行处理能力强、运算速度高、可在室温工作、能耗仅为电子计算机的几十分之一，但目前尚未诞生。
（3）量子计算机：利用量子叠加效应和量子并行计算，数据用量子位存储，一个量子位可同时存储 0 和 1，存储量比普通计算机大许多，运算速度可能比个人计算机快 10 亿倍；存在退相干和纠错两大困难，但已见曙光。
（4）纳米计算机：用纳米技术研制芯片，体积仅数百个原子大小，几乎不需要耗能，性能比今天的计算机强大许多，且芯片生产成本低。
（2）重点发展领域（4 分）：教材 2.6.2 列出光子计算机、智能及神经网络计算机、模糊逻辑和多值处理计算机等方向，并指出计算机未来发展要与多种科学领域结合，重点是：高性能计算器件；高性能计算机的体系结构（并行计算、集群系统）；高性能的软件；和谐的人机交互环境（智能接口，如计算机视觉、语音识别、虚拟现实）；大力发展纳米技术；网络技术与多媒体性能。
（3）对农业信息化的影响（4 分）：教材 2.6 指出计算机将由信息处理、数据处理过渡到知识处理，知识库将取代数据库，自然语言、模式、图像、手写体等人机会话将成为输入输出的主要形式。据此，新型计算技术将从多方面支撑农业信息化：高性能计算与并行计算可支撑农业大数据分析、作物生长模型模拟和大范围气象预报；量子计算有望加速种植方案优化、物流路径规划等组合优化问题；生物计算与纳米计算可大幅降低计算与传感设备的能耗和成本，适合田间部署；光子计算可提升遥感图像处理与传输能力；和谐人机交互（语音、图像识别）可降低农户使用信息系统的门槛，推动农业信息服务向智能化、个性化发展。
\begin{analysis}得分点：（1）四个发展方向及各自核心思想、优势 4 分；（2）重点发展领域答出四至六项 4 分；（3）结合农业应用分析 4 分。常见失分：把“主要发展方向”（生物、光子、量子、纳米）与“重点发展领域”（高性能器件、体系结构、软件、人机交互、纳米技术、网络多媒体）混为一谈；只写技术名词而无优势或影响分析。本题对应 2021 年真题“计算机技术的重点发展领域”。教材 2.6，第 25—29 页。\end{analysis}

\end{document}
